\section{Protótipo}
\label{sec:prototipo}

\subsection{Visão geral}

O protótipo tem três partes. A \textbf{biblioteca} (C++23) implementa as camadas
como elementos de uma cadeia de processamento: cada mensagem que chega passa pela
limitação de tráfego, pela camada de segurança (o envelope) e, opcionalmente, pela
persistência; cada camada pode ser ligada e desligada em execução, sem reconstruir
a cadeia. O \textbf{\emph{gateway} de ensaio} liga essa cadeia a um \emph{broker}
Mosquitto real através de um cliente Paho e reporta, por janela de medição, as
mensagens entregues e rejeitadas (por motivo), a latência e o tempo de processador.
O \textbf{laboratório de carga} (Python) lança o \emph{broker}, o \emph{gateway},
os sensores simulados e o atacante, aplica cada configuração e mede.

A Figura~\ref{fig:arquitetura} mostra onde corre cada parte na campanha.

\begin{figure}[ht]
  \centering
  \begin{tikzpicture}[
      proc/.style={draw=cinza, fill=white, rounded corners=2pt, font=\footnotesize,
                   align=center, minimum height=0.85cm, inner sep=2pt},
      camada/.style={draw=iscte, fill=fundotab, rounded corners=2pt, font=\footnotesize,
                     align=center, text width=4.1cm, minimum height=0.62cm, inner sep=2pt},
      maquina/.style={draw=iscte, thick, rounded corners=4pt, inner sep=6pt},
      rotulo/.style={font=\small\bfseries, text=iscte, anchor=south west},
      seta/.style={-{Stealth[length=2mm]}, thick, iscte},
      dupla/.style={{Stealth[length=2mm]}-{Stealth[length=2mm]}, thick, iscte},
      controlo/.style={-{Stealth[length=1.6mm]}, densely dashed, cinza},
      nota/.style={font=\scriptsize, text=cinza, align=center}]
    % --- computador de secretária
    \node[proc, text width=3.1cm] (s1) at (1.75, 0.9)
      {Agente de sensores 1\\\scriptsize 25 sensores e atuadores};
    \node[proc, text width=3.1cm] (s2) at (1.75, -0.4)
      {Agente de sensores 2\\\scriptsize 25 sensores e atuadores};
    \node[nota, below=0.08cm of s2] (env) {envelope em Python};
    \node[proc, text width=3.1cm, draw=mau] (atk) at (1.75, -2.3)
      {Agente atacante\\\scriptsize estranho ou dispositivo\\comprometido};
    \node[maquina, fit=(s1)(s2)(env)(atk)] (pc) {};
    \node[rotulo] at (pc.north west) {Computador de secretária};
    % --- Raspberry Pi 5
    \node[proc, text width=2.1cm, fill=iscte!10] (brk) at (6.25, 0.9)
      {\emph{Broker}\\Mosquitto 2.0};
    \node[nota, below=0.08cm of brk, text width=2.4cm] (brkn)
      {TLS/mTLS, autenticação, ACL, limites};
    \node[proc, text width=2.0cm, draw=mau, dashed] (rel) at (8.95, 0.9)
      {Intermediário\\\scriptsize só em ataques};
    \node[font=\scriptsize\bfseries, text=cinza] (gwt) at (12.85, 1.75)
      {\emph{Gateway} de ensaio (C++23)};
    \node[camada, fill=white, draw=cinza] (paho) at (12.85, 0.9) {Cliente MQTT (Paho)};
    \node[camada, below=0.2cm of paho] (lim) {Limitação de tráfego};
    \node[camada, below=0.2cm of lim] (sec) {\textbf{Segurança da mensagem}\\
      \scriptsize envelope, validade, janela,\\ políticas de modo e de chave};
    \node[camada, below=0.2cm of sec] (per) {Persistência (opcional)};
    \node[camada, fill=white, draw=cinza, below=0.2cm of per] (med)
      {Aplicação de ensaio\\\scriptsize entregas, rejeições,\\ latência, processador};
    \node[maquina, draw=cinza, thin, inner sep=4pt, fit=(gwt)(paho)(med)] (gw) {};
    \node[proc, text width=4.6cm] (ctl) at (7.6, -2.3)
      {Controlador do laboratório\\\scriptsize Python; configura, mede e guarda};
    \node[maquina, fit=(brk)(brkn)(rel)(gw)(ctl)] (pi) {};
    \node[rotulo] at (pi.north west) {Raspberry~Pi~5};
    % --- fluxos MQTT (medidos)
    \draw[dupla] (s1.east) -- ([yshift=0.2cm]brk.west);
    \draw[dupla] (s2.east) -- ([yshift=-0.05cm]brk.west);
    \draw[seta, mau] (atk.east) -- ([yshift=-0.3cm]brk.west);
    \draw[dupla] (brk) -- (rel);
    \draw[dupla] (rel) -- (paho);
    \foreach \a/\b in {paho/lim, lim/sec, sec/per, per/med}
      \draw[seta] (\a) -- (\b);
    % --- controlo do laboratório (fora do caminho medido)
    \draw[controlo] (ctl.west) -- (ctl.west -| pc.east);
    \draw[controlo] (ctl.north) -- (ctl.north |- brkn.south);
    \draw[controlo] (ctl.east) -- (ctl.east -| gw.west);
  \end{tikzpicture}
  \caption{Arquitetura implementada, tal como corre na campanha. A cadeia de camadas
    da Parte~1 (\PUm{Secção~6.1}) está no \emph{gateway} de ensaio. Setas
    contínuas: MQTT (telemetria dos sensores para o \emph{gateway}, ordens no
    sentido inverso; a vermelho, o atacante). O intermediário só existe nos ensaios
    de ataques na rede. A tracejado: o canal de controlo do laboratório (TLS), que
    lança e configura cada parte e fica fora do caminho medido.}
  \label{fig:arquitetura}
\end{figure}

\subsection{O que foi implementado}

A Tabela~\ref{tab:implementado} relaciona cada elemento da proposta com a sua
concretização.

\begin{table}[ht]
  \centering
  \caption{Elementos da arquitetura proposta e a sua implementação.}
  \label{tab:implementado}
  \footnotesize
  \begin{tabularx}{\textwidth}{@{}L{4.6cm}YL{4.2cm}@{}}
    \toprule
    \textbf{Elemento (Parte~1)} & \textbf{Implementação} & \textbf{Onde} \\
    \midrule
    Envelope de 25 octetos; tópico e modo nos dados associados; modos de cifra,
      autenticação e assinatura (§6.4) & AES-256-GCM, GMAC e Ed25519 sobre OpenSSL;
      a mesma implementação em Python para os sensores, verificada pelos mesmos
      vetores de teste & \texttt{modules/security}, \texttt{msec\_envelope.py} \\
    ChaCha20-Poly1305 como alternativa (§6.4) & O algoritmo é um atributo da chave:
      o formato do envelope não muda e um recetor serve os dois & \texttt{chacha20\_poly1305\_}\allowbreak\texttt{cipher.cpp} \\
    Ordem de verificação (§6.5) & Formato, modo aceite, validade, verificação
      criptográfica, chave--tópico, janela & \texttt{security\_layer.cpp} \\
    Validade temporal (§6.6) & $\Delta_a$ e $\Delta_s$ globais e por fluxo (filtro
      de tópico) & \texttt{FreshnessRule} \\
    Janela anti-repetição, $w = 64$ (§6.6) & Por chave, em memória, com estado
      limitado & \texttt{replay\_guard.cpp} \\
    Restrição chave--tópico (§6.7) & Política com padrões por dono da chave,
      aplicada depois da autenticação & \texttt{key\_policy.cpp} \\
    Chaves por época, rotação e revogação (§6.8) & Rotação com período de graça;
      revogação de chaves, palavras-passe e certificados (CRL) & \texttt{pki.py}, \texttt{lab.py} \\
    Limitação de tráfego antes da criptografia (§6.9) & \emph{Token bucket} por
      tópico & \texttt{modules/limits} \\
    Camada de transporte (§6.3) & TLS~1.3, autenticação mútua, palavras-passe,
      ACL por utilizador, identificador de cliente ligado ao utilizador, limites de
      pacote e de ligações & \texttt{broker.py} (Mosquitto) \\
    \bottomrule
  \end{tabularx}
\end{table}

\subsection{Decisões tomadas nesta parte}

\paragraph{Reinício do recetor.} A Parte~1 registou como limitação que a janela
anti-repetição se perde quando o recetor reinicia. A solução adotada não persiste a
janela a cada mensagem: o recetor guarda de forma durável apenas a época de chaves
mais recente que aceitou e, depois de um reinício, só aceita épocas posteriores. A
gravação faz-se uma vez por rotação, antes de aceitar a primeira mensagem da nova
época, e uma falha de leitura ou de escrita rejeita a mensagem.

\paragraph{Dependências entre mecanismos.} A validade temporal, a anti-repetição
e a restrição chave--tópico verificam campos do envelope. Uma configuração que as
ligue com o envelope desligado é recusada, no arranque e em execução, em vez de as
ignorar em silêncio.

\paragraph{Fora do âmbito.} Não foram implementados um modo HMAC próprio (o GMAC
cobre a autenticação sem cifra), \emph{nonces} por contador nem o modo de cifra
seguida de assinatura, que a Parte~1 apresenta como possível extensão.

\subsection{Testes automatizados}

O protótipo tem 93 testes em C++
(unitários, de integração e de aceitação), executados com AddressSanitizer e
UndefinedBehaviorSanitizer, e 73 testes em Python (61 do laboratório, 12 da
análise). Os testes de aceitação têm um
teste por regra da ordem de verificação e foram verificados por mutação: cada um
falha quando a regra correspondente é removida. Três alvos de \emph{fuzzing}
exercitam a análise do envelope, a validação de tópicos e todo o caminho de
receção com octetos arbitrários. A integração contínua compila com
\texttt{clang-tidy} (avisos tratados como erros) e corre os cenários de ataque
contra um Mosquitto real.
